\documentclass[11pt]{article}
\usepackage[margin=1in]{geometry}
\usepackage{amsmath}
\usepackage{amssymb}
\usepackage{enumitem}
\usepackage{graphicx}
\usepackage{xcolor}
\usepackage{tcolorbox}

\title{Problem 1b: The Expenditure Function \\ (Starting from Scratch)}
\author{ECO 310 - For People Who Skipped Lecture}
\date{February 23, 2026}

\begin{document}

\maketitle

\section{What Are We Doing Again?}

Part (a) asked: "What should you buy to reach utility $u^*$ cheaply?"

Part (b) asks: "Okay, now how much does that cost?"

That's it! The expenditure function is just the cost of what you found in part (a).

\section{Super Basic Definition}

\begin{tcolorbox}[colback=blue!5!white,colframe=blue!75!black,title=Expenditure Function]
\textbf{Expenditure function} = How much money you need to spend to reach utility level $u^*$

It's literally: (price of good 1) $\times$ (amount of good 1) + (price of good 2) $\times$ (amount of good 2)
\end{tcolorbox}

\textbf{Formula:}
$$e(p_1, p_2, u^*) = p_1 \cdot x_1^c + p_2 \cdot x_2^c$$

where $x_1^c$ and $x_2^c$ are from your answer to part (a).

\section{Starting with Examples (Coffee and Tea)}

Let's use actual numbers first, then write the general formula.

\subsection{Example 1: Coffee is Cheaper}

\textbf{Given:}
\begin{itemize}
\item Coffee costs \$3 per cup ($p_1 = 3$)
\item Tea costs \$5 per cup ($p_2 = 5$)
\item You want utility 10 ($u^* = 10$)
\item Utility function: $u = \text{coffee} + \text{tea}$
\end{itemize}

\textbf{From Part (a):} Since coffee is cheaper ($3 < 5$), you buy:
$$(x_1^c, x_2^c) = (10, 0)$$
That means: 10 coffees, 0 teas.

\textbf{Part (b) - Calculate expenditure:}
\begin{align*}
e(3, 5, 10) &= (\text{price of coffee}) \times (\text{coffees bought}) + (\text{price of tea}) \times (\text{teas bought}) \\
&= 3 \times 10 + 5 \times 0 \\
&= 30 + 0 \\
&= \$30
\end{align*}

\textbf{Answer:} It costs \$30 to reach utility level 10.

\subsection{Example 2: Tea is Cheaper}

\textbf{Given:}
\begin{itemize}
\item Coffee costs \$5 per cup ($p_1 = 5$)
\item Tea costs \$2 per cup ($p_2 = 2$)
\item You want utility 10 ($u^* = 10$)
\end{itemize}

\textbf{From Part (a):} Since tea is cheaper ($2 < 5$), you buy:
$$(x_1^c, x_2^c) = (0, 10)$$
That means: 0 coffees, 10 teas.

\textbf{Part (b) - Calculate expenditure:}
\begin{align*}
e(5, 2, 10) &= 5 \times 0 + 2 \times 10 \\
&= 0 + 20 \\
&= \$20
\end{align*}

\textbf{Answer:} It costs \$20 to reach utility level 10.

\subsection{Example 3: Different Utility Target}

\textbf{Given:}
\begin{itemize}
\item Coffee costs \$3 per cup ($p_1 = 3$)
\item Tea costs \$5 per cup ($p_2 = 5$)
\item You want utility 25 ($u^* = 25$)
\end{itemize}

\textbf{From Part (a):} Since coffee is cheaper, you buy:
$$(x_1^c, x_2^c) = (25, 0)$$

\textbf{Part (b) - Calculate expenditure:}
\begin{align*}
e(3, 5, 25) &= 3 \times 25 + 5 \times 0 \\
&= 75 + 0 \\
&= \$75
\end{align*}

\textbf{Answer:} It costs \$75 to reach utility level 25.

\section{See the Pattern?}

Look at the examples:
\begin{itemize}
\item When coffee is cheaper (\$3 vs \$5): cost = (\$3) $\times$ (utility target)
\item When tea is cheaper (\$2 vs \$5): cost = (\$2) $\times$ (utility target)
\end{itemize}

\textbf{Pattern:} cost = (cheaper price) $\times$ (utility target)

\section{Writing the General Formula}

Now let's do this without specific numbers. Use $p_1$, $p_2$, and $u^*$ as variables.

\subsection{Recall Part (a) Answer}

From part (a), compensated demand is:

$$(x_1^c, x_2^c) = \begin{cases}
(u^*, 0) & \text{if } p_1 < p_2 \\
(0, u^*) & \text{if } p_1 > p_2
\end{cases}$$

\subsection{Calculate Expenditure for Each Case}

The expenditure function is: $e(p_1, p_2, u^*) = p_1 \cdot x_1^c + p_2 \cdot x_2^c$

\textbf{Case 1: If $p_1 < p_2$ (good 1 is cheaper)}

From part (a): $(x_1^c, x_2^c) = (u^*, 0)$

Plug into expenditure formula:
\begin{align*}
e(p_1, p_2, u^*) &= p_1 \cdot u^* + p_2 \cdot 0 \\
&= p_1 u^*
\end{align*}

\textbf{Case 2: If $p_1 > p_2$ (good 2 is cheaper)}

From part (a): $(x_1^c, x_2^c) = (0, u^*)$

Plug into expenditure formula:
\begin{align*}
e(p_1, p_2, u^*) &= p_1 \cdot 0 + p_2 \cdot u^* \\
&= p_2 u^*
\end{align*}

\textbf{Case 3: If $p_1 = p_2$ (same price)}

From part (a): any combination works, like $(u^*, 0)$ or $(0, u^*)$ or $(u^*/2, u^*/2)$

But they all cost the same:
\begin{align*}
e(p_1, p_2, u^*) &= p_1 u^* = p_2 u^* \quad \text{(since $p_1 = p_2$)}
\end{align*}

\section{Final Answer to Part (b)}

\begin{tcolorbox}[colback=yellow!10!white,colframe=orange!75!black,title=Answer: Expenditure Function (Piecewise)]
\begin{equation}
e(p_1, p_2, u^*) = \begin{cases}
p_1 u^* & \text{if } p_1 < p_2 \\
p_2 u^* & \text{if } p_1 > p_2 \\
p_1 u^* = p_2 u^* & \text{if } p_1 = p_2
\end{cases}
\end{equation}
\end{tcolorbox}

\subsection{Even Simpler Form}

You can write this more compactly as:

\begin{tcolorbox}[colback=yellow!10!white,colframe=orange!75!black,title=Answer: Expenditure Function (Compact)]
\begin{equation}
e(p_1, p_2, u^*) = \min\{p_1, p_2\} \cdot u^*
\end{equation}
\end{tcolorbox}

\textbf{Translation:} Minimum cost = (whichever price is lower) $\times$ (target utility)

Both forms are correct! The compact form is cleaner.

\section{Verify with Our Examples}

Let's check that our formula works:

\subsection{Example 1 Check}

Coffee \$3, tea \$5, $u^* = 10$

Using formula: $e(3, 5, 10) = \min\{3, 5\} \times 10 = 3 \times 10 = 30$ \checkmark

\subsection{Example 2 Check}

Coffee \$5, tea \$2, $u^* = 10$

Using formula: $e(5, 2, 10) = \min\{5, 2\} \times 10 = 2 \times 10 = 20$ \checkmark

\subsection{Example 3 Check}

Coffee \$3, tea \$5, $u^* = 25$

Using formula: $e(3, 5, 25) = \min\{3, 5\} \times 25 = 3 \times 25 = 75$ \checkmark

Perfect! Our formula works.

\section{Why Does This Make Sense?}

\begin{tcolorbox}[colback=green!5!white,colframe=green!75!black,title=Intuition]
\textbf{Question:} What's the cheapest way to get $u^*$ units of utility?

\textbf{Answer:} With perfect substitutes, buy only the cheaper good.

\textbf{Cost:} (cheaper price) $\times$ (how many units you need) = (cheaper price) $\times$ $u^*$
\end{tcolorbox}

This is exactly what $e(p_1, p_2, u^*) = \min\{p_1, p_2\} \cdot u^*$ says!

\section{Step-by-Step Process for Part (b)}

\begin{enumerate}
\item Look at your answer from part (a) - what bundle do you buy in each case?
\item For each case, calculate total cost: $p_1 x_1 + p_2 x_2$
\item Simplify the expressions
\item Write as piecewise function OR use $\min\{p_1, p_2\}$ notation
\item Done!
\end{enumerate}

\section{Common Mistakes to Avoid}

\begin{itemize}
\item \textbf{Don't} forget to multiply by $u^*$ - the cost depends on how much utility you want!
\item \textbf{Don't} average the prices - you buy only the cheaper good, not a mix
\item \textbf{Don't} use $m$ (income) - expenditure function uses $u^*$ (target utility), not budget
\end{itemize}

\section{Summary}

\begin{tcolorbox}[colback=blue!5!white,colframe=blue!75!black,title=TL;DR]
\textbf{Part (a):} Find what to buy $\rightarrow$ $(u^*, 0)$ if good 1 cheaper, $(0, u^*)$ if good 2 cheaper

\textbf{Part (b):} Calculate the cost $\rightarrow$ $e(p_1, p_2, u^*) = \min\{p_1, p_2\} \cdot u^*$

\textbf{Intuition:} Buy the cheaper good, so cost = (cheaper price) $\times$ (target utility)
\end{tcolorbox}

\end{document}
